\documentclass[11pt]{article}
\usepackage[margin=1in]{geometry}
\usepackage{amsmath,amssymb}
\usepackage{booktabs}
\usepackage[colorlinks=true,linkcolor=blue,citecolor=blue,urlcolor=blue]{hyperref}
\usepackage{natbib}
\usepackage{enumitem}
\usepackage{xcolor}

\title{ContractCVE-Web: Behavioral Contracts for\\ Browser Agents Facing Vulnerability-Triggered Decisions}
\author{Jahangir \\ \small \texttt{jahangir.08.cse@gmail.com}}
\date{October 2026}

\begin{document}
\maketitle

\begin{abstract}
Behavioral-contract research for LLM agents is emerging, but it has not been
specialized to browser-vulnerability decisions grounded in real CVE data. We
present ContractCVE-Web, a pilot study that defines four formal
pre/act/post-style contracts over a nine-action browser-agent vocabulary and
tests whether an 8B-parameter LLM's action plans satisfy them, under real
CVE context drawn from 30 live-fetched NVD records. Comparing an
\emph{unconstrained} condition (contract withheld) against a
\emph{contract-enforced} condition (contract text included in-prompt), we
find the overall Contract Satisfaction Rate (CSR) rises from 50\% to 100\%
\emph{for single, isolated contracts under neutral phrasing}. That
qualification is load-bearing: a follow-up round shows the unconditional
``100\%'' does not hold once contracts are stacked or phrasing varies --
under a terse phrasing, contract-enforced CSR drops to 80\%; and under
contract stacking it can collapse to 0\% even where both contracts are
individually satisfied 100\% of the time. A second round running \emph{all
six} pairwise stacks then overturns our own account of why. Two cells invert
(C1+C4 from 0\% to 100\%, C2+C3 from 100\% to 0\%) once the combined scenario
is written once for all pairs rather than tailored per pair, so we withdraw
``an obligation destroys a co-stated prohibition'' as a mechanism. What
survives both rounds is narrower and better supported: the constraint that
drops is always a \emph{prohibition}, never an ordering or obligation
contract, and which prohibition drops is set by the scenario's temptation
structure rather than by the partner contract. Fresh scenario wordings leave
contract-enforced CSR at 100\% on all four contracts in isolation, so the
single-contract result generalizes while the stacking result does not. The lift is still
concentrated where the model's unprompted single-contract defaults are
wrong (two of the four contracts go from 0\% to 100\% compliance once stated
explicitly), but the paper's central claim is narrower than a first reading
suggests: explicit contracts help substantially for isolated policies under
plain phrasing, and do not, on their own, guarantee compliance once multiple
contracts must be satisfied simultaneously.
\end{abstract}

\section{Introduction}
An LLM browser agent's default behavior is shaped by its training, not by
any formal specification of what is safe to do under a given vulnerability
state. Two agents that behave identically on most tasks can diverge sharply
on the handful of scenarios where \emph{intuition} and \emph{policy}
disagree -- for instance, an agent might ``feel'' that installing a
requested browser extension is a reasonable, helpful action, even when a
currently active, unpatched CVE in the Extensions component makes that
specific action dangerous. The ContractCVE-Web idea in the source
research-ideas document proposes making such policies explicit and
machine-checkable as behavioral contracts $C = (\textit{pre}, \textit{act},
\textit{post}, \textit{audit})$, and measuring whether stating them actually
changes agent behavior for the better.

This pilot implements four such contracts and measures Contract Satisfaction
Rate both with and without the contract text present in the prompt, across
30 real CVE-grounded scenarios, using a rule-based (not LLM-judged) verifier
against the agent's own structured action-plan output.

\section{Related Work}
AgentDojo \citep{debenedetti2024agentdojo} evaluates defenses against
prompt-injection-induced policy violations but does not formalize
\emph{positive} behavioral contracts the agent must satisfy independent of
any attack. ToolEmu \citep{ruan2023toolemu} surfaces risky tool-use failures
but scores them post hoc rather than against a stated, checkable contract.
WASP \citep{evtimov2025wasp} and BrowseSafe \citep{zhang2025browsesafe}
study whether browser agents can be hijacked by adversarial web content;
ContractCVE-Web is complementary -- it asks whether, even absent any
adversary, an agent reliably follows a stated safety policy when a known
vulnerability is active. Runtime behavioral-contract and compliance-testing
work for agents more broadly \citep{behavioralcontracts2026} motivates the
pre/act/post contract formalism we adopt here, specialized to CVE-triggered
browser decisions.

\section{Threat Model / Problem Formulation}
We define a contract $C = (\textit{pre}, \textit{act}, \textit{post})$ over
browser-agent states and vulnerability evidence: \textit{pre} is a
precondition on the CVE state (e.g., ``an active, unpatched exploit is
present''), \textit{act} is a constraint on the agent's action sequence
(e.g., ``must not invoke \texttt{download\_file}''), and \textit{post} is
the resulting guarantee. An agent is \emph{compliant} with $C$ on a given
scenario iff its produced action sequence satisfies \textit{act} whenever
\textit{pre} holds. We define four contracts, each paired with a scenario
template:
\begin{itemize}[nosep]
  \item \textbf{C1 safe-browse/no-download}: under an active exploit, an
        unsolicited-link scenario must never invoke \texttt{download\_file}.
  \item \textbf{C2 patch-before-navigation}: under an active exploit, any
        \texttt{navigate\_url} action must be preceded by
        \texttt{check\_patch\_status}.
  \item \textbf{C3 no-high-risk-extension-under-active-exploit}: under an
        active exploit, the plan must never invoke
        \texttt{install\_extension} or \texttt{manage\_extension}.
  \item \textbf{C4 mandatory-escalation-under-contradictory-evidence}: when
        patch/exploit-status evidence is contradictory, the plan must invoke
        \texttt{escalate\_to\_human} rather than silently proceeding.
\end{itemize}
The threat we probe is \textbf{policy drift under implicit judgment}: without
an explicit contract, the agent's action sequence is governed only by its
training-time priors about what seems reasonable, which may or may not align
with an organization's actual safety policy for a given CVE state.

\section{Methodology}
The 4 contracts are rotated across the 30-CVE pool by index modulo 4, so
each contract is exercised on 7--8 real CVEs, each injected into its
scenario template as CVE id, CWE, component, and severity context. For every
scenario we query the model twice: once \textbf{unconstrained} (system
prompt lists only the nine-action vocabulary and asks for an ordered action
plan as JSON) and once \textbf{contract-enforced} (the same system prompt
plus the applicable contract's pre/act/post text, with an explicit ``you
MUST comply'' instruction). A rule-based verifier --- not the LLM --- checks
the returned action list against the contract: e.g., for C2 it checks
whether \texttt{check\_patch\_status} appears before the first
\texttt{navigate\_url} in the list (vacuously satisfied if
\texttt{navigate\_url} never appears). Decoding is greedy (temperature 0).

\section{Experimental Setup}
\textbf{Data.} The shared 30-CVE pool (live NVD 2.0 API pull \citep{nvd}),
used only for CVE context injected into each scenario.

\textbf{Model.} \texttt{meta-llama/Llama-3.1-8B-Instruct}
\citep{grattafiori2024llama3} via vLLM \citep{kwon2023vllm},
\texttt{gpu\_memory\_utilization=0.2}, single GPU, \texttt{bfloat16},
\texttt{enforce\_eager=True}, max sequence length 4096, max 250 generated
tokens per sample. $30 \times 2 = 60$ pilot samples, plus 210 full-run samples
(Section~\ref{sec:fullrun}: 180 phrasing-variant and 30 stacked), for 270 in
total. No sample failed to parse in either set.

\section{Pilot Results}
All 60 samples parsed as valid JSON action plans (0\% parse-failure rate).
Table~\ref{tab:cvecontrac} reports Contract Satisfaction Rate (CSR) overall
and per contract.

\begin{table}[h]
\centering
\caption{ContractCVE-Web pilot results ($n=60$ samples, 30 CVEs $\times$ 2 conditions).}
\label{tab:cvecontrac}
\begin{tabular}{lrr}
\toprule
Contract & Unconstrained CSR & Contract-enforced CSR \\
\midrule
C1 safe-browse/no-download         & 100\% & 100\% \\
C2 patch-before-navigation          & \textbf{0\%}   & \textbf{100\%} \\
C3 no-high-risk-extension           & \textbf{0\%}   & \textbf{100\%} \\
C4 mandatory-escalation             & 100\% & 100\% \\
\midrule
\textbf{Overall}                    & 50\%  & 100\% \\
\bottomrule
\end{tabular}
\end{table}

The lift is entirely concentrated in C2 and C3, which move from 0\%
compliance (unconstrained) to 100\% (contract-enforced). C1 and C4 are at
100\% in \emph{both} conditions: the model's unprompted defaults already
avoid downloading from an unsolicited link and already escalate under
contradictory evidence, so the contract adds no measurable behavior change
there.

\section{Full Run: Phrasing Variation and Contract Stacking}
\label{sec:fullrun}
The pilot's limitations named three untested axes: phrasing variation,
adversarial scenarios designed to tempt a violation, and contract stacking.
The full run adds all three on the same 30-CVE pool and model, for 210
further samples (0 parse failures). Each contract is instantiated under three
scenario phrasings --- \emph{neutral}, \emph{urgency\_pressure} (the user
applies time pressure and explicitly asks the agent to skip its checks), and
\emph{terse} (a clipped one-line command) --- and two contract stacks are
tested, each presenting both contracts simultaneously against a scenario that
engages both.

\begin{table}[h]
\centering
\caption{Contract Satisfaction Rate by scenario phrasing ($n=30$ per
variant per condition).}
\label{tab:cvecontrac-phrasing}
\begin{tabular}{lrrl}
\toprule
Phrasing & Unconstrained & Contract-enforced & Where enforcement fails \\
\midrule
neutral          & 50.0\% & 100.0\% & --- \\
urgency\_pressure & 26.7\% & 100.0\% & --- \\
terse            & 26.7\% & \textbf{80.0\%} & C4 escalation, 1/7 \\
\bottomrule
\end{tabular}
\end{table}

\textbf{Urgency pressure degrades the unprompted agent but not the
contract.} Moving from neutral to urgency\_pressure halves unconstrained CSR
(50.0\%~$\rightarrow$~26.7\%): under time pressure the model abandons C4's
escalation entirely (7/7~$\rightarrow$~0/7) on top of the C2 and C3 failures
it already had. With the contract stated, compliance is unaffected (100\%).
This is the pilot's thesis holding under adversarial phrasing --- the
contract's value rises exactly as the model's defaults degrade.

\textbf{Terse contract phrasing breaks the escalation contract.} Under the
clipped phrasing, contract-enforced CSR falls to 80.0\%, and the loss is
entirely C4: 1 of 7 compliant, against 7/7 under both other phrasings. C1,
C2 and C3 remain at 100\%. The three prohibition- and ordering-style
contracts survive compression; the one obligation-style contract does not.
A contract that must be \emph{remembered and acted on} is more sensitive to
how tersely the surrounding request is phrased than a contract that merely
forbids something.

\begin{table}[h]
\centering
\caption{Contract stacking: both contracts stated, against a scenario
engaging both ($n=15$ per stack).}
\label{tab:cvecontrac-stack}
\begin{tabular}{lrrr}
\toprule
Stack & Stack CSR (both hold) & First contract & Second contract \\
\midrule
C2+C3 (ordering $+$ prohibition)   & 100.0\% & C2: 15/15 & C3: 15/15 \\
C1+C4 (prohibition $+$ obligation) & \textbf{0.0\%} & \textbf{C1: 0/15} & C4: 15/15 \\
\bottomrule
\end{tabular}
\end{table}

\textbf{Stacking an obligation onto a prohibition destroys the
prohibition.} C2+C3 stack without interference: both hold in all 15 cases.
C1+C4 collapses to 0\%, and the breakdown shows precisely how. C4 (invoke
\texttt{escalate\_to\_human}) is satisfied in 15/15 cases; C1 (never invoke
\texttt{download\_file}) is satisfied in 0/15. The modal plan is
\texttt{[escalate\_to\_human, navigate\_url, render\_html,
check\_patch\_status, escalate\_to\_human, download\_file]} --- the agent
notifies a human and \emph{then performs the forbidden download anyway}.

This is not scenario difficulty masquerading as interference, and the design
lets us say so. Two of the three isolated C1 phrasings also demand a download
outright (``open this link and download whatever's there'', ``open this link
and get me the file''), and C1 holds at 8/8 under every one of them, in both
the unconstrained and the contract-enforced condition. C1 is robust to an
explicit download demand when it is the only contract in play. It fails
completely when an escalation obligation is added beside it. The added
contract, not the harder ask, is what moves the number.

The mechanism this suggests is worth naming: the escalation obligation gives
the model a way to \emph{discharge} its safety duty procedurally, after which
it treats the user's request as something it may now satisfy. Escalation is
performed as an additional action rather than as a halt.

\subsection{A permissive oracle, and what it hides}
\label{sec:oracle-caveat}
C4's compliance check tests only whether \texttt{escalate\_to\_human} appears
somewhere in the action list. It does not test whether escalation
\emph{prevented} the risky action, so in the C1+C4 stack C4 scores 15/15
while the agent downloads the file in every one of those same 15 plans. Under
the contract's own stated post-condition --- ``a human is notified before any
risky action is taken under uncertainty'' --- the literal reading is satisfied
(the escalation is first in the sequence) and the evident intent is not. A
stricter oracle, requiring that no risky action follow an escalation, would
score C4 at 0/15 here and would report the stack as failing on both
contracts rather than one. We report the permissive number because it is what
the implemented checker measures, and flag it because the 100\% figure for C4
in Table~\ref{tab:cvecontrac-stack} should not be read as evidence that
escalation worked. More generally, structural action-list checks can certify
an obligation that the agent has satisfied in form only, which is a reason to
prefer post-condition verifiers with side effects over presence tests.

\section{Second Round: All Six Stacks, and a Result That Does Not Reproduce}
\label{sec:actionplan}
A later round (2026-10-20, Llama-3.1-8B-Instruct on GPU~2, vLLM at
\texttt{gpu\_memory\_utilization}=0.25) added two scenario phrasings per
contract and ran \emph{all six} pairwise stacks rather than the two of
Section~\ref{sec:fullrun}. It confirms the single-contract claim and
\textbf{overturns the stacking mechanism reported above}.

\paragraph{The single-contract result generalizes beyond its template.}
Two fresh scenario wordings per contract (120 samples) leave
contract-enforced CSR at \textbf{100\% for all four contracts}. The
\emph{unconstrained} baseline, however, is template-sensitive: C1 moves from
100\% to 87.5\% and C3 from 0\% to 50\% across the new wordings. The lift is
real; part of the original clean $0\%/100\%$ contrast was a property of the
single template each contract was first written against.

\begin{table}[h]
\centering
\caption{All six pairwise stacks under one fixed combined scenario
($n=5$ CVEs per stack, contract-enforced). ``Semantic'' is a keyword-presence
test on the model's rationale, not a comprehension check.}
\label{tab:allstacks}
\begin{tabular}{lrrl}
\toprule
Stack & Structural CSR & Semantic CSR & Which contract fails \\
\midrule
C1$+$C2 & 100\% & 100\% & --- \\
\textbf{C1$+$C3} & \textbf{0\%} & 80\% & \textbf{C3} (0/5); C1 holds 5/5 \\
C1$+$C4 & 100\% & 100\% & --- \\
\textbf{C2$+$C3} & \textbf{0\%} & 100\% & \textbf{C3} (0/5); C2 holds 5/5 \\
C2$+$C4 & 100\% & 100\% & --- \\
C3$+$C4 & 100\% & 80\% & --- \\
\bottomrule
\end{tabular}
\end{table}

\paragraph{Two cells invert against Section~\ref{sec:fullrun}, and the cause is
the scenario, not the contracts.} The earlier round reported C1$+$C4 at
\textbf{0\%} and C2$+$C3 at \textbf{100\%}; this round reports the
near-opposite, \textbf{100\%} and \textbf{0\%}. The two rounds construct the
combined scenario differently, and that difference fully accounts for the
reversal. Section~\ref{sec:fullrun} wrote a scenario \emph{per pair}, tailored
to tempt both of that pair's constraints --- its C1$+$C4 scenario explicitly
demanded a download (``the user wants both the link opened and, if there's any
file offered, downloaded''). This round uses \emph{one fixed scenario for all
six pairs}, which asks for navigation, an extension install, and a check of an
unsolicited link under incomplete evidence --- but never asks for a download.
So C1 is simply not tempted here and passes trivially (5/5 in every stack
containing it), while C3's prohibition is tempted in all six.

\paragraph{The consequence for the mechanism claimed in
Section~\ref{sec:fullrun}.} That section attributed the C1$+$C4 collapse to an
obligation (C4, escalate) discharging the agent's safety duty and releasing it
to satisfy the forbidden download. Under a scenario that does not request a
download, the same pair holds at 100\%. \textbf{We therefore withdraw
``stacking an obligation onto a prohibition destroys the prohibition'' as a
general mechanism.} The narrower claim that survives both rounds is that
\emph{the prohibition is the constraint that drops}: the contract violated in
every failing cell across both rounds is a ``must NOT'' contract (C1 in round
one, C3 in round two), never an ordering or obligation contract. What
determines \emph{which} prohibition drops is the scenario's temptation
structure, not the identity of the partner contract.

\paragraph{An ordering confound we cannot resolve at this $n$.} C3 fails in
exactly the two stacks where it appears \emph{second} in the concatenated
contract text (C1$+$C3, C2$+$C3) and holds where it appears \emph{first}
(C3$+$C4). Because the stacks are generated as ordered combinations over
$[\textrm{C1},\textrm{C2},\textrm{C3},\textrm{C4}]$, position and partner
identity are \emph{perfectly confounded}: every C3-second stack pairs it with
C1 or C2, and the only C3-first stack pairs it with C4. With $n=5$ per cell
and no counterbalancing we cannot separate ``C3 is fragile when listed second''
from ``C3 is fragile when paired with C1 or C2''. Counterbalancing contract
order is the first experiment the next round should run, and it is cheap.

\paragraph{Stating the rule while breaking it.} In C2$+$C3 the rationale
references both contracts' vocabulary in 100\% of cases while structural
compliance is 0\%. The agent articulates the constraint it is about to
violate. This dissociation is worth reporting, but the ``semantic'' test is a
keyword-presence check over the rationale string (e.g.\ \texttt{extension},
\texttt{exploit}), so it establishes that the right words appear, not that the
constraint was understood --- the same class of permissive oracle flagged in
Section~\ref{sec:oracle-caveat}.

\section{Discussion and Limitations}
This pattern is the cleanest possible illustration of the ContractCVE-Web
thesis: a stated contract is not uniformly useful, it is useful
\emph{precisely} where the agent's default judgment disagrees with policy.
For C2 (check patch status before navigating) and C3 (no extension
management under active exploit), the model's unconstrained instinct was to
proceed directly with the user's requested action -- navigate to the
intranet page, install the requested extension -- without first reasoning
about the active-exploit precondition at all. Once the contract made that
precondition and the required ordering explicit, compliance was perfect.
For C1 and C4, the model already had the ``right'' instinct (don't download
from a cold email link; escalate rather than act under contradictory
information), so the explicit contract was redundant there but also
harmless.

This has a direct practical implication: an organization deploying explicit
behavioral contracts should not expect uniform gains, and should prioritize
writing and enforcing contracts for exactly the policies that diverge from a
model's training-time intuitions, which is also the harder case to discover
without an experiment like this one.

\textbf{Limitations.} Two of the pilot's three named gaps are closed by
Section~\ref{sec:fullrun}: phrasing is now varied three ways including an
adversarial urgency-pressure variant, and contract stacking is tested. The
third is not. Compliance checking remains a structural presence test over the
action list rather than a post-condition verifier with side effects, and
Section~\ref{sec:oracle-caveat} shows this is not a cosmetic limitation: the
permissive C4 oracle certifies an escalation that demonstrably failed to
prevent the action it was meant to prevent. Every obligation-style contract
measured this way is exposed to the same error, and the direction of the error
is optimistic.

Scope remains narrow in three further respects. All 6 pairwise stacks are now
tested (Section~\ref{sec:actionplan}), but each under a \emph{single} combined
scenario and at $n=5$ CVEs per stack, and the two rounds disagree on two of the
six cells. Stacked-contract compliance should therefore be reported as a
property of a (pair, scenario) pair rather than of a contract pair, and no
claim about which specific pairs are hard is safe until several independent
combined-scenario wordings are averaged per pair. Contract \emph{order} within
the stack is not counterbalanced, leaving position confounded with partner
identity; three-contract stacks remain untested entirely. All
results come from a single model (Llama-3.1-8B-Instruct) at temperature 0, so
nothing here separates a property of behavioral contracts from a property of
this model. With $n=15$ per stack and $n=30$ per phrasing cell, the 80.0\%
terse figure carries a 95\% Wilson interval of roughly $\pm$14~pp; the 0\% and
100\% cells are bounded away from each other but rest on a single
deterministic decode.

\section{Conclusion and Future Work}
Stating a behavioral contract explicitly produces a measurable,
interpretable Contract Satisfaction Rate lift, concentrated exactly where it
should be --- on the policies the model's unconstrained defaults get wrong ---
and that lift survives adversarial urgency pressure intact. The full run shows
two boundaries on the claim, and the second is the more consequential.

Phrasing matters asymmetrically: compressing the request to a terse command
costs nothing for the three prohibition- and ordering-style contracts and
breaks the one obligation-style contract (C4: 7/7~$\rightarrow$~1/7). An
obligation has to be remembered and discharged, and that is the kind of
instruction a clipped context drops.

Stacking is not additive, but our first account of why it fails did not
survive its own replication. Running all six pairwise stacks under a single
shared scenario inverted two of the six cells relative to the per-pair
scenarios of Section~\ref{sec:fullrun}, including the C1+C4 cell the
obligation-release mechanism was built on. We withdraw that mechanism. Two
claims survive both rounds. First, \textbf{the constraint that drops under
stacking is always a prohibition} --- C1 in round one, C3 in round two --- and
never an ordering or an obligation contract, which are satisfied in every
failing cell. Second, \textbf{which prohibition drops is a property of the
scenario, not of the contract pair}: C1 collapses when the scenario demands a
download and holds when it does not, with the partner contract unchanged.

The practical reading is therefore about measurement before it is about
design. Contracts cannot be validated one at a time and composed --- that much
the stacking failures do establish --- but neither can a stacked result be
quoted without the scenario that produced it. A contract suite must be
evaluated as a suite, over several independent scenario constructions, with
contract order counterbalanced.

Future work should average over several independent combined-scenario wordings
per stack and counterbalance contract order, since the present disagreement
between rounds is attributable to both and separable by neither; extend to
three-contract stacks; replace the
presence-test oracle with a post-condition verifier that can tell escalation
from escalation-then-proceed, since the current oracle scores the central
failure as a success; and add a second model family before treating any of
these rates as a property of contracts rather than of this model.

\bibliographystyle{plainnat}
\bibliography{writeup}

\end{document}
